%% Technical appendix / supplementary material for SpecBraM.
%% Compile as its own PDF; submit separately from the 7-page main paper.
\documentclass[letterpaper]{article}
\usepackage[submission]{aaai2027}
\usepackage[hyphens]{url}
\usepackage{graphicx}
\urlstyle{rm}
\def\UrlFont{\rm}
\usepackage{natbib}
\usepackage{caption}
\frenchspacing
\usepackage{booktabs}
\usepackage{multirow}
\usepackage{amsmath,amssymb}
\usepackage{xcolor}
\newcommand{\best}[1]{\textbf{#1}}
\newcommand{\val}[2]{#1\,{\scriptsize$\pm$#2}}
\newcommand{\snd}[1]{\underline{#1}}
\newcommand{\na}{--}
\pdfinfo{/TemplateVersion (2027.1)}

\title{Technical Appendix\\SpecBraM: Predicting Band Power, Not Codes or Waveforms\\A Data- and
Label-Efficient Spectral-Prediction Foundation Model for Clinical EEG}
\author{Anonymous Submission}
\affiliations{Paper under double-blind review}

\begin{document}
\maketitle
\appendix

\section{Datasets}\label{app:datasets}
\noindent\textbf{Pretraining corpus.} We pretrain on the Temple University Hospital
EEG Corpus (TUEG, v2.0.1) \cite{obeid2016tuh}. Every recording is mapped to the
$19$-channel $10$--$20$ montage, resampled to $256$\,Hz, and cut into
non-overlapping $10$\,s windows. Our primary model uses a random $2{,}100$-patient
subset ($2{,}388$ hours; $859{,}606$ windows) drawn with a fixed seed; the
full-data control uses the full corpus ($14{,}764$ patients,
$\sim17{,}666$ hours). All four downstream benchmarks are separate corpora that are
disjoint from TUEG at the patient level, so no benchmark leakage is possible for any
TUEG-pretrained method.

\noindent\textbf{Downstream benchmarks.} Table~\ref{tab:app_data} summarizes the
four leak-free benchmarks; we adopt the subject-disjoint train/validation/test
splits of CBraMod \cite{wang2025cbramod} and CSBrain \cite{csbrain}, described below.

\noindent\emph{ISRUC-Sleep (subgroup~I)} \cite{khalighi2016isruc}: overnight
polysomnography from $100$ subjects scored into five stages (W, N1, N2, N3, REM) on
$30$\,s epochs. The six mastoid-referenced EEG derivations (F3/C3/O1--A2,
F4/C4/O2--A1) are placed by name into our $19$-channel $10$--$20$ layout and the
remaining channels are zero-padded; each $30$\,s epoch is cropped to its central
$10$\,s to match the pretraining window. We use the subject-disjoint split
(train $1$--$80$, val $81$--$90$, test $91$--$100$) and evaluate
sequence-to-sequence.

\noindent\emph{HMC} \cite{alvarez2021hmc}: the Haaglanden Medisch Centrum sleep
database, $151$ whole-night recordings with the same five-stage, $30$\,s scoring. We
select the four EEG derivations (F4--M1, C4--M1, O2--M1, C3--M2) by channel name,
zero-pad the rest, and discard EMG/EOG/ECG; epochs are again cropped to their
central $10$\,s. The split is $100$/$25$/$26$ (train/val/test), matching CSBrain.


\noindent\emph{Mental Arithmetic (EEGMAT)} \cite{zyma2019mentalarith}: $36$ subjects
recorded at rest and during serial-subtraction mental arithmetic (binary
rest-vs.-arithmetic). Its discriminative signal is spectral (alpha suppression and
frontal-midline theta), placing it in scope. The native $19$-channel $10$--$20$
recordings ($500$\,Hz) map cleanly to our montage after resampling to $256$\,Hz; we
use $5$\,s windows and the CSBrain split (subjects $1$--$28$/$29$--$32$/$33$--$36$
for train/val/test).

\begin{table*}[t]\centering\footnotesize
\setlength{\tabcolsep}{6pt}
\caption{Downstream benchmarks; test-subject counts in parentheses. ``seq2seq'' =
sequence-to-sequence sleep staging over $30$\,s epochs.}
\label{tab:app_data}
\begin{tabular}{llccl}
\toprule
Dataset & Task & Cls. & Subjects & Window\,/\,runs\\
\midrule
ISRUC-I \cite{khalighi2016isruc}          & Sleep staging & 5 & 100 (10) & 30\,s seq2seq / 3 reps\\
HMC \cite{alvarez2021hmc}                 & Sleep staging & 5 & 151 (26) & 30\,s seq2seq / 3 reps\\
Ment.\ Arith.\ \cite{zyma2019mentalarith} & Rest vs.\ arith.\ & 2 & 36 (4) & 5\,s / 5 seeds\\
\bottomrule
\end{tabular}
\end{table*}

\section{Preprocessing}\label{app:preproc}
All corpora pass through one uniform pipeline. Channels are mapped into a canonical
$19$-channel $10$--$20$ order \emph{by name} (never by position), missing channels
are zero-padded, and non-EEG channels (EMG/EOG/ECG) are discarded; this name-based
matching avoids the positional-slicing errors that can silently mix in non-EEG or
mis-ordered channels. Signals are resampled to $256$\,Hz and normalized per
recording with a robust scaler, $x\mapsto(x-\mathrm{median})/(\mathrm{IQR}/1.349)$,
computed on the \emph{full} recording before windowing; robust (median/IQR) scaling
preserves the long-range amplitude structure that per-window $z$-scoring destroys.
Artifact filtering follows CBraMod: recordings shorter than $5$\,min are dropped,
$60$\,s are trimmed from each recording's start and end (to remove
montage/settling artifacts), and windows whose samples exceed $\pm150\,\mu$V are
rejected; no notch filter is applied.

\noindent\textbf{Window length.} The tokenizer emits one token per $0.78$\,s patch,
and the encoder's learned positional embeddings admit up to $64$ time patches, so the
model accepts variable-length windows without any retraining: a $10$\,s window
becomes $12$ patches and a $5$\,s window $6$. We therefore keep each benchmark's
native window---Mental Arithmetic $5$\,s---and crop only the $30$\,s
sleep epochs to their central $10$\,s, following the CBraMod/CSBrain sleep-staging
protocol. No window is padded or stretched to a fixed $10$\,s length.

\section{Architecture and Pretraining Hyperparameters}\label{app:hp}
\noindent\textbf{Filterbank tokenizer.} Each channel is convolved with $N{=}5$
learnable depthwise \texttt{Conv1d} band-pass filters (kernel $65$), initialized to
Hann-windowed sinusoids centred at $2/6/10/21/38$\,Hz (the $\delta\theta\alpha\beta
\gamma$ bands) and then trained freely, so the model may sharpen or shift the bands
to fit the data. Each of the five band signals is cut into $12$ non-overlapping
$200$-sample ($0.78$\,s) patches, linearly embedded to $d_b{=}64$, and the five
per-band embeddings of a patch are concatenated and projected (with LayerNorm) to
the $D{=}512$ token, yielding a $[19\times12\times512]$ token grid.

\noindent\textbf{Criss-cross encoder.} Twelve transformer blocks each apply
self-attention along the channel axis and then along the time axis, rather than full
$19{\times}12$ attention; this models spatial (montage) and temporal structure at
$O(C^2T+CT^2)$ instead of $O((CT)^2)$ cost, with learned positional embeddings on
both axes. Masked time patches replace their band embeddings with a shared learnable
mask token before encoding. A single linear head maps each encoded token to five
predicted log band-powers, supervised only on masked patches against a \emph{fixed}
(frozen, no-gradient) reference filterbank applied to the raw signal; this head is
discarded after pretraining and only the encoder transfers downstream.

\noindent\textbf{Hyperparameters.} Table~\ref{tab:app_hp} lists all architecture and
pretraining hyperparameters. The objective is
$\mathcal{L}=\mathcal{L}_{\mathrm{sp}}+0.05\,\mathcal{L}_{\mathrm{sig}}$, with no
JEPA or reconstruction term. Pretraining runs for $105$ scheduled epochs, but the
downstream model is the epoch-$16$ checkpoint at which pretraining loss converges;
later checkpoints overfit the small subset (Sec.~4 of the main paper). This single
checkpoint, selected by pretraining loss (blind to downstream test performance), is
used for all tasks.

\begin{table}[!ht]\centering\footnotesize
\setlength{\tabcolsep}{4pt}
\caption{Architecture and pretraining hyperparameters (SpecBraM).}
\label{tab:app_hp}
\begin{tabular}{ll}
\toprule
\multicolumn{2}{l}{\textit{Architecture}}\\
\midrule
Learnable filterbank        & $N{=}5$ depthwise Conv1d, kernel $65$\\
Band centres (init)         & $2/6/10/21/38$\,Hz ($\delta\theta\alpha\beta\gamma$)\\
Per-band embedding $d_b$    & $64$\\
Token dimension $D$         & $512$\\
Encoder                     & $12$ criss-cross layers, $8$ heads\\
Input                       & $10$\,s, $256$\,Hz, $19$ ch ($2560$)\\
Patch $P$ / \#patches $T_p$ & $200$ ($0.78$\,s) / $12$\\
Parameters                  & $51.2$\,M\\
\midrule
\multicolumn{2}{l}{\textit{Objective}}\\
\midrule
Loss                        & $\mathcal{L}_{\mathrm{sp}}+0.05\,\mathcal{L}_{\mathrm{sig}}$\\
Mask axis / ratio           & time patches / $0.30$\\
Band-power weight           & $1.0$\\
SIGReg $\lambda$            & $0.05$\\
JEPA / reconstruction       & none\\
\midrule
\multicolumn{2}{l}{\textit{Optimization}}\\
\midrule
Optimizer                   & AdamW\\
Learning rate               & $1.7{\times}10^{-3}$\\
Weight decay                & $0.05$\\
Warmup / schedule           & $7$ epochs / cosine\\
Gradient clipping           & $3.0$\\
Batch (per-GPU / effective) & $64$ / $512$\\
Epochs (trained / reported) & $105$ / $16$\\
Hardware                    & $8\times$ RTX A6000 (48\,GB)\\
Seed                        & $42$\\
\bottomrule
\end{tabular}
\end{table}

\section{Fine-Tuning and Evaluation Protocols}\label{app:ft}
For every downstream task we report both a \emph{frozen probe} and \emph{full
fine-tuning}, and against a \emph{random-initialized} baseline that uses the
identical architecture and protocol (isolating the contribution of pretraining).

\noindent\textbf{Frozen probe.} We freeze the encoder and train only a linear/head
classifier on its features. This is deterministic (no backbone stochasticity) and
serves as our primary readout on the tiny clinical cohorts, where full fine-tuning
is high-variance.

\noindent\textbf{Full fine-tuning.} The whole network is updated with AdamW; the
per-benchmark schedules are in Table~\ref{tab:app_ft}. Sleep (ISRUC/HMC) uses
linear-probe-then-fine-tune: the head is first trained on frozen features, then the
encoder is unfrozen and the model is optimized under a cosine schedule (base
learning rate $10^{-4}$, batch size $8$). Mental Arithmetic follows the LaBraM-style protocol with
layer-wise learning-rate decay ($0.65$) and a $5$-epoch warmup. All variants use up
to $50$ epochs and select the checkpoint with the best validation balanced accuracy.
Because Mental Arithmetic ($4$ test subjects) is
tiny and thus high-variance, we run $5$ seeds and report mean$\pm$std, whereas the
larger sleep sets use $3$ reps.

\noindent\textbf{Metrics.} We report balanced accuracy (mean per-class recall) as
the primary metric, since the clinical classes are imbalanced; for the five-class
sleep tasks we additionally report Cohen's $\kappa$ and the sample-weighted F1, and
for the binary tasks the areas under the precision-recall and ROC curves (AUC-PR,
AUROC). Sleep staging is evaluated sequence-to-sequence over full nights.

\begin{table}[t]\centering\scriptsize
\setlength{\tabcolsep}{3pt}
\caption{Fine-tuning protocol per benchmark. Protocols and subject-disjoint splits
follow CBraMod/CSBrain \citep{wang2025cbramod,csbrain}, so SpecBraM and the reported
baselines are fine-tuned under the same recipe. ``decay'' $=$ layer-wise
learning-rate decay.}
\label{tab:app_ft}
\begin{tabular}{lcc}
\toprule
                & Sleep (ISRUC/HMC)  & Ment.\ Arith.\\
\midrule
Protocol        & LP\,$\to$\,full FT  & LaBraM-style\\
Optimizer       & AdamW  & AdamW\\
Base LR         & $1{\times}10^{-4}$ & $5{\times}10^{-4}$\\
Schedule        & cosine  & decay $0.65$\\
OneCycle max    & -- &  --\\
Warmup epochs   & $0$  & $5$\\
Max ep.\ / patience & $50$ / -- & $50$ / $50$\\
Batch size      & $8$ & $64$\\
Selection       & best-val BA &  best-val BA\\
Window          & $30$\,s seq2seq  & $5$\,s\\
Runs            & $3$ reps & $5$ seeds\\
\bottomrule
\end{tabular}
\end{table}

\section{Code and Reproducibility}\label{app:code}
All code for pretraining, preprocessing, and every downstream evaluation is provided
as (anonymized) supplementary material and will be released publicly under the MIT
license upon publication, with inline comments referencing the equations and
sections of this paper. Experiments ran on $8\times$NVIDIA RTX A6000 ($48$\,GB) GPUs,
an Intel Xeon Gold 5220R CPU, Ubuntu 22.04 (Linux 5.15), PyTorch 2.12.0 with
CUDA 13.0, and MNE-Python 1.12.1. Random seeds (base $42$, with per-run offsets for
the multi-seed evaluations) are set for NumPy, PyTorch, and CUDA.

\bibliography{reference}

\end{document}
